\documentclass[11pt]{article}
\usepackage[margin=1.05in]{geometry}
\usepackage{amsmath,amsthm,amssymb,mathtools,bm}
\usepackage{booktabs,array,enumitem}\usepackage[expansion=false]{microtype}
\usepackage{xcolor}
\usepackage[colorlinks=true,linkcolor=blue!60!black,citecolor=blue!60!black,urlcolor=blue!60!black]{hyperref}
\theoremstyle{plain}
\newtheorem{theorem}{Theorem}[section]\newtheorem{proposition}[theorem]{Proposition}\newtheorem{lemma}[theorem]{Lemma}
\newtheorem{corollary}[theorem]{Corollary}\newtheorem{conjecture}[theorem]{Conjecture}
\theoremstyle{definition}\newtheorem{definition}[theorem]{Definition}\newtheorem{remark}[theorem]{Remark}
\newtheorem{example}[theorem]{Example}
\newcommand{\R}{\mathbb R}\newcommand{\Z}{\mathbb Z}\newcommand{\E}{\mathbb E}\newcommand{\Var}{\operatorname{Var}}\newcommand{\Cov}{\operatorname{Cov}}
\newcommand{\relu}{\operatorname{relu}}\newcommand{\diag}{\operatorname{diag}}\newcommand{\tr}{\operatorname{tr}}
\newcommand{\Tr}{\operatorname{Tr}}\newcommand{\1}{\mathbf 1}\newcommand{\cA}{\mathcal A}\newcommand{\cB}{\mathcal B}
\newcommand{\cF}{\mathcal F}\newcommand{\cT}{\mathcal T}\newcommand{\cH}{\mathcal H}\newcommand{\aff}{\mathfrak{aff}}
\newcommand{\He}{\mathrm{He}}\newcommand{\Fa}{\Phi_{\mathrm F}}\newcommand{\cZ}{\mathcal Z}\newcommand{\Sym}{\operatorname{Sym}}
\title{\textsc{Localization in the commutant, stage 13:}\\ \large the leaderboard system read through convex geometry.\\
The Schur hub as Eldan's cumulant drift, the critical gain as a conserved dilation charge,\\ and the network as a difference of
support functions}
\author{}
\date{10 October 2026}
\begin{document}
\maketitle
\begin{abstract}
We do three things. First, we take apart the current leaderboard system (v56) and separate what it inherits, what we engineered,
and which pieces realize theory. Second, we read six works in full: Klartag--Lehec on the slicing problem, Bizeul--Klartag--Lehec on
KLS, Klartag--Ordentlich on the strong data-processing inequality under heat flow, Klartag's stochastically evolving ellipsoid,
Sahasrabudhe's survey of exponentially small scales, and the caustic-collar paper on billiard rigidity. Third, we prove or check
the bridges they open into our setting.
\begin{itemize}[leftmargin=*]
\item \emph{The Schur hub is Eldan's cumulant drift.} Along Eldan's localization, the drift of the $m$-th cumulant is
$-(m\kappa_m+L_m)$, where $L_m$ contracts pairs of lower cumulants through the inverse covariance (Bizeul--Klartag--Lehec).
We show this is the law of total cumulance taken infinitesimally along the localization path. For $m=4$, $L_4$ is three
$\kappa_3A^{-1}\kappa_3$ pairings; checked to $10^{-7}$ on a discrete measure, where it is $20\%$ of the drift. This is the
algebraic form of v56's largest single gain, the Schur-hub path class $12\,[YC^{-1}Y^\top]_{ii}$.
\begin{itemize}
\item It explains three measured facts: why the hub needs $C^{-1}$, why it needs the source-resolved (quenched) third
cumulant rather than the annealed one, and why its derived weight is right.
\item It predicts the missing pairings for the $(2,2)$ and $(3,1)$ slices, with their coefficients.
\item It says $\kappa_3$ has no such term ($L_3=0$), so the chain's third-cumulant errors are transport errors, as measured.
\end{itemize}
In the two-tier language of stage 12, the hub is the between-tier fourth cumulant: v56's biggest gain is the two-tier
framework realized at fourth order.
\item \emph{The critical gain is a conserved dilation charge.} \textsc{ReLU} commutes with scale mixing, so the network's common
random scale is transported exactly. That is why the mean-direction mode is critical (transport factor $\approx0.95$ measured).
\begin{itemize}
\item Stage 12's increment law gives its size: across-input log-norm increments $\frac4n(\frac32-\frac{J_2(\theta)}{2\pi})\approx4\theta^2/n$, which
saturate. This answers the open item of note XLIV with a parameter-free leading law: $0.0058$--$0.0072$ at layers $8$--$14$,
against measured $0.0077$--$0.0100$; the input radius gives $1/(2n)$.
\item The gain need not mean-revert to explain the plateau. A falsifiable prediction separates the two readings.
\item The Edgeworth-corrected Gaussian closure is exact to $O(\varepsilon^2)$ on a scale mixture when it carries the mixture's own
skewness $6\varepsilon r$ and kurtosis $12\varepsilon$ (checked: halving $\varepsilon$ cuts the error by $3.8$--$4.0$).
\end{itemize}
So the critical mode leaks into the means only through inconsistency between the chain's carried variance, $\kappa_3$ and
$\kappa_4$ along the dilation. A dilation-consistent (scale-mixture) closure removes the leak at $O(n^2)$ per layer.
\item \emph{The spherical geometry is a support-function geometry.} By homogeneity the problem lives on the sphere.
\begin{itemize}
\item A bias-free \textsc{ReLU} network is a difference of support functions of two polytopes (its tropical Newton polytopes),
and its Gaussian mean is a difference of mean widths, i.e.\ first intrinsic volumes. Stage 11's border formula is the edge
formula for this intrinsic volume.
\item Each layer evaluates the support function of the zonoid of the previous layer's law at the next rows. The Gaussian
closure replaces that zonoid by a Gaussian (ellipsoidal) one.
\item The curvature of the zonoid is the wall density of stages 11--12. All three identities are checked exactly.
\end{itemize}
Klartag's evolving ellipsoid and the Lovett--Meka walk share one principle: once the process touches a constraint, it keeps it
and continues in the constraint-preserving subspace. That principle is the right treatment of the chain's exact invariants (the
dilation charge, the Euler--Stein ladder).
\item \emph{What the billiard paper transfers} is the support-function machinery, the grazing parameter $\lambda=t^2$ (our
temperature $T\approx\theta^2/2$), and the method of pinning high modes by family invariants. It brings no new algorithm by itself.
\item \emph{Heat flow and slicing.} We record a computable memory budget, the negentropy as an integrated localization
deficit, checked to $0.5\%$.
\end{itemize}
Each claim is labelled as proved, known, measured, conjectural or an analogy. The checks are in \texttt{scripts/verify\_stage13.py}.
\end{abstract}
\tableofcontents
\section{The picture in one page}
\label{sec:summary}
\paragraph{1. What v56 is (Section~\ref{sec:v56}).}
\begin{itemize}[leftmargin=*]
\item \textbf{The base.} A public chain (504aldo's v29): Gaussian moment closure, $\kappa_3$ sources as CP legs in young and old tiers,
hub slices, memoryless $\kappa_4$, a Wick table and Strassen-for-billing. We added cost engineering.
\item \textbf{Five theory-derived terms on top.} The quenched $\kappa_4$ pair class, the scale-mixture ``pair'', the fold (exact
first-order $(2,1)$ feedback as an arm displacement), $103$ renormalization counterterms, and the chaos-graded $\kappa_4$
diagonal (the Schur hub plus the third-chaos star).
\item \textbf{Where the gain came from.} The two large wins ($-10\%$ each, held out) are exact derived terms whose refitted amplitudes stay
at $0.95$ and $1.007$. Shapes chosen by convention never paid.
\item \textbf{What is not yet inside.} Stages 10--12's quantitative structure is not inside, and nothing yet treats the critical
dilation mode.
\end{itemize}

\paragraph{2. The Schur hub is Eldan's cumulant drift (Section~\ref{sec:hierarchy}).} Along Eldan's normalized localization,
$d\kappa_m=\kappa_{m+1}(\cdot,A^{-1/2}dB)-(m\kappa_m+L_m)dt$, and we show $L_m$ is the between-path part of the law of total cumulance.
\begin{itemize}[leftmargin=*]
\item $L_3=0$, so $\kappa_3$'s errors are transport errors.
\item $L_4$ is three $\kappa_3A^{-1}\kappa_3$ pairings, three times the variance of the regression of the square on the field. This is the
form of v56's path class $12[YC^{-1}Y^\top]_{ii}$, and it explains why the hub needs $C^{-1}$, the quenched (source-resolved) $\kappa_3$, and
weight $1$.
\item It predicts the $(2,2)$ and $(3,1)$ pairings, with coefficients, from the same solve.
\item In the two-tier language, v56's largest gain is the between-tier fourth cumulant.
\end{itemize}

\paragraph{3. The critical gain is a conserved dilation charge (Section~\ref{sec:dilation}).} \textsc{ReLU} commutes with scale
mixing, so the common scale is transported with factor exactly $1$.
\begin{itemize}[leftmargin=*]
\item Its size follows stage 12's increment law: $\frac4n(\frac32-\frac{J_2(\theta)}{2\pi})\approx4\theta^2/n$ per layer, saturating at $nV\to44.6$; the
input radius contributes $1/(2n)$.
\item Predicted values are $0.0058$--$0.0072$ at layers $8$--$14$, against measured $0.0070$--$0.0100$.
\item The plateau needs no mean reversion; the test is stated.
\item The Edgeworth-corrected closure is exact to $O(\varepsilon^2)$ on a scale mixture carrying its own $\gamma_1=6\varepsilon r$ and
$\gamma_2=12\varepsilon$. So the charge moves the means only through inconsistency of the carried variance, $\kappa_3$ and $\kappa_4$.
\item Design rule: impose the scale-mixture package, at $O(n^2)$ per layer. The truth is about $70\%$ package, with kurtosis about
$1.7\times$ the skewness slope against the package's $2$.
\end{itemize}

\paragraph{4. The spherical geometry (Section~\ref{sec:support}).}
\begin{itemize}[leftmargin=*]
\item The network is $h_{P_+}-h_{P_-}$, a difference of support functions of tropical polytopes, and $\E f=(V_1(P_+)-V_1(P_-))/\sqrt{2\pi}$, a
difference of first intrinsic volumes. Stage 11's border formula is the edge formula of $V_1$.
\item Each layer evaluates the support function of its law's zonoid at the next rows. The closure is the Gaussian (ellipsoidal)
zonoid, and the zonoid's curvature is the wall density.
\item Klartag's evolving ellipsoid gives a dictionary (contacts as frozen gates, free directions as unresolved gates) and one
principle: keep every exact constraint you touch.
\end{itemize}

\paragraph{5. Billiards, packing, heat flow (Sections~\ref{sec:billiards}--\ref{sec:heat}).}
\begin{itemize}[leftmargin=*]
\item The caustic paper's support-function envelopes, its positivity $g+g''>0$ and its $\lambda=t^2$ are ours: the wall density, and the
temperature $T\approx\theta^2/2$.
\item Its rigidity has a probabilistic twin: a closure exact on a full family of tilts forces Gaussianity.
\item Its Newton-on-high-modes method needs family invariants we lack. That is the flex theorem from the other side; v56 supplied the
missing data through source-resolved legs.
\item From slicing: the negentropy is an integrated localization deficit, checked to $0.5\%$, a candidate memory budget.
\item From Shannon--Stam stability: the quenched content is the variance of the localized covariance, whose diffusion coefficient
is $\kappa_3$.
\end{itemize}

\paragraph{6. What to build (Section~\ref{sec:next}).}
\begin{enumerate}[leftmargin=*]
\item The dilation-consistent closure.
\item The remaining $L_4$ pairings from v56's existing solve.
\item Sticky projection onto the exact identities.
\item The memory-budget diagnostic.
\end{enumerate}

\section{What v56 is: inherited, engineered, and theory-realized parts}
\label{sec:v56}
The system is the baked file \texttt{estimator\_\allowbreak final\_\allowbreak v56.py} in \texttt{notes/\allowbreak leaderboard-\allowbreak system} on branch
\texttt{claude/\allowbreak bold-\allowbreak albattani-\allowbreak xgecxi}, made from \texttt{est\_v29.py} with one configuration string. Scored on all $100$ official networks it has adjusted MSE
$3.137\times10^{-9}$ in-sample for its counterterms, and $3.166\times10^{-9}$ by the held-out protocol (raw $1.56\times10^{-8}$,
$C/B=0.203$). That is $-24\%$ against our V35 and roughly $-42\%$ against the public v29 it started from. The public leaders are
quoted at $1.1$--$2.1\times10^{-9}$ (note XXXIX).

\subsection{What it computes, layer by layer}
The state carried from layer to layer is a truncated cumulant description of the pre-activations $z_l$:
\begin{itemize}[leftmargin=*]
\item the mean $\mu_l$ and the full covariance $C_l$;
\item third-cumulant \emph{sources}: each layer gives birth to a $\kappa_3$ ``star'' $\Sym(C_b\otimes I\otimes C_b\diag(w_2))$, the Price/Hermite
term of \textsc{ReLU} at the Gaussian point, carried as CP-factor legs $(A_s,P_s)$;
\item a few slices of the fourth cumulant (a diagonal and some $(2,2)$ and $(3,1)$ shapes).
\end{itemize}
One layer then runs these steps:
\begin{enumerate}[leftmargin=*]
\item \textbf{Linear step.} $\mu\leftarrow W\mu$ and $C\leftarrow WCW^\top$ exactly. The legs are transported by the mean-gate propagator
$B=W\diag(\Phi)$, which is first-order linear response, the ``product gate''. Young sources stay dense. Old sources are compressed to a
shared rank-$320$/$192$ Lyapunov basis, re-orthogonalized at one randomized range-finder \emph{join} per layer.
\item \textbf{Slices.} The $(3)$ and $(2,1)$ slices of $\kappa_3$ (D3, D21) are read from the legs by Hadamard ``hub'' products.
\item \textbf{Wick stage.} The post-activation moments come from a Gaussian-plus-Edgeworth closure in the carried cumulants: a
Hermite/$\sigma$ table, Mehler order $2$, giving K11, K21, K22, K1--K4.
\item \textbf{Birth.} A new source is born from the closure. A rank-$4$ residual of the $(2,1)$ slice is kept, and $\kappa_4$ is fed back
into $\kappa_3$.
\item \textbf{Fourth-order sector.} $\kappa_4$ is regenerated without memory: annealed transport of the $(2,2)$ content plus an
adaptively fitted $\lambda$, with corrections.
\end{enumerate}
All large products are written as Strassen recursions so that flopscope bills their analytic count.

\subsection{Inherited, engineered, and theory-realized}
\begin{center}\small
\begin{tabular}{>{\raggedright\arraybackslash}p{4.3cm}>{\raggedright\arraybackslash}p{6.0cm}>{\raggedright\arraybackslash}p{4.3cm}}\toprule
component & origin & measured worth\\\midrule
\multicolumn{3}{l}{\emph{Inherited from the public chain (504aldo v29; probably shared with other leaders)}}\\
Gaussian moment closure with Hermite/Wick tables & standard (the organisers' cumulant-propagation reference) & the backbone\\
$\kappa_3$ sources as transported CP legs, young/old tiers, range-finder joins & 504aldo v29 & the backbone\\
D3/D21 hub slices; memoryless $\kappa_4$ with fitted $\lambda$; birth residual; feedback & 504aldo v29 & the backbone\\
Strassen-for-billing & 504aldo v29 & the cost model\\\midrule
\multicolumn{3}{l}{\emph{Engineered by us (cost and numerics, theory-light)}}\\
warm-started joins (Lyapunov continuity), Strassen level 6, symmetric/aliased Grams, billing fixes, dead-neuron saturation drop,
leaf minima, join schedule, retuned ranks & notes XV, XXIX, V29--V35 & $-3$ to $-6\%$ adjusted from warm joins and level 6; $C/B$ from $0.267$ (v29) to $0.180$ (V35)\\\midrule
\multicolumn{3}{l}{\emph{Theory realized in code}}\\
quenched $\kappa_4$ pair class (\texttt{V33\_K4Q}) & retained transport, note XXVIII: quenched over annealed & $-2.3\%$\\
the ``pair'' (\texttt{V33\_WK4M=3}, \texttt{V31\_K4D=3}) & gain-aware / scale-mixture shapes, notes XI, XIX, XXXVIII & $-0.9\%$\\
the fold (\texttt{V52\_FB\_FOLD=1}) & exact first-order $(2,1)$ feedback as a displacement of the arm, $a\to a+Y_t/2w_2+X_t/6$, with the arm's double duty & $-10.0\%$ held-out adjusted\\
counterterms (\texttt{V47\_CAL}, $103$ amplitudes) & renormalization: fitted elements as counterterms of the closure & $-3.8\%$\\
chaos-graded $\kappa_4$ diagonal (\texttt{V56\_P4}) & Wiener chaos grading: Schur-hub path class plus third-chaos star at derived weights & $-9.8\%$ held-out adjusted ($-20.9\%$ raw)\\\bottomrule
\end{tabular}
\end{center}
Three remarks.
\begin{enumerate}[leftmargin=*]
\item \textbf{The two large wins are exact, derived terms whose fitted amplitude stays at its derived value.} The fold refits at $0.95$
and the chaos correction at $1.007$. In note XXXIX's sense they are law terms. The conventions (shapes of omitted classes chosen from
the pair state) never paid, as the flex theorem predicted.
\item \textbf{Where the programme's ideas are already inside.}
\begin{itemize}
\item The quenched/annealed distinction of stage 9 is realized twice. The quenched $\kappa_4$ pair class wins, and in the chaos
correction the source-resolved hub product wins while its symmetrized (annealed) control carries nothing.
\item The quasi-free (Fock/Bogoliubov) sector of stages 11--12 is realized as the chaos grading.
\item Section~\ref{sec:hierarchy} shows that the Schur hub is the between-tier fourth cumulant of the two-tier (localization)
decomposition.
\end{itemize}
\item \textbf{What is not inside yet.}
\begin{itemize}
\item Nothing of stages 10--12's quantitative structure: the Mehler clock, the parabolic depth, the internal Eldan clock and
temperature, the hyperbolic upper tier, the wall-density law.
\item In particular, nothing yet treats the critical dilation mode that note XLIV measured. Section~\ref{sec:dilation} supplies its law
and a closure that respects it.
\end{itemize}
\end{enumerate}

\section{What the six works do}
\label{sec:readings}
\paragraph{Klartag--Lehec, the slicing problem.}
\begin{itemize}[leftmargin=*]
\item \textbf{Localization setup.} For an isotropic log-concave $\mu$, localization tilts by $e^{\theta\cdot x-t|x|^2/2}$ with
$d\theta_t=dW_t+a_tdt$. The covariance process $A_t$ obeys $\E A_t+\E a_t\otimes a_t=\mathrm{Id}$ (the variance decomposition) and
$\frac d{dt}\E A_t=-\E A_t^2$.
\item \textbf{The final step.} Guan's bound $\E\Tr A_t^2\le Cn$ plus a projection-monotonicity lemma produce a subspace and a
time $t_0\sim1$ with $\E A_{t_0}\ge\frac14\mathrm{Id}$, without losing the isotropic constant (Milman's $M$-ellipsoids).
Eldan--Mikulincer's stability of Shannon--Stam then bounds $\int\E|\Gamma_r-\E\Gamma_r|^2/(1-r)\,dr$, the fluctuation of the localized
covariance, by the entropy gain under $X\mapsto(X+X')/\sqrt2$, which Ball--Nguyen bound by $2n$.
\item \textbf{The identity we use} (their Corollary 3.2 with Lemma 3.3): the relative entropy to the Gaussian is the integrated
F\"ollmer drift, which in localization time reads
\[
D(\mu\,\|\,\gamma)=\tfrac12\int_0^\infty\Big(\tfrac n{1+t}-\E\Tr A_t\Big)dt\qquad(\mu\ \text{isotropic}).
\]
\end{itemize}

\paragraph{Bizeul--Klartag--Lehec, KLS.}
\begin{itemize}[leftmargin=*]
\item \textbf{The base case.} The third-cumulant bound $R_3\le8\,\mathrm{Id}$ comes from the sharp generalized thin shell.
\item \textbf{The induction.} It is propagated to $R_m\le K^{m-1}((m-1)!)^2\mathrm{Id}$ through Eldan's normalized localization
$dc=A^{-1}a\,dt+A^{-1/2}dB$, $dQ=A^{-1}dt$. Its key \emph{cumulant equation} (their Lemma 4.3) is
\begin{gather*}
d\kappa_m=\kappa_{m+1}(\cdot,\dots,\cdot,A_t^{-1/2}dB_t)-\big(m\kappa_m+L_m\big)dt,\\
L_m(h)=\!\!\!\sum_{\substack{J\ni1\\2\le|J|\le m-2}}\!\!\!\big\langle\kappa_{|J|+1}(h_J,\cdot),A_t^{-1}\kappa_{m-|J|+1}(h_{J^c},\cdot)\big\rangle .
\end{gather*}
\item \textbf{The finish.} A suspension turns cumulant bounds into bounds on the Taylor tensors of tilt averages
$F_f(z)=\int fe^{\langle z,x\rangle}d\mu/\int e^{\langle z,x\rangle}d\mu$. The Song--Zhang criterion then gives a dimension-free Poincar\'e
constant.
\item \textbf{The Gaussian case.} For the Gaussian input, $F_f(z)=\E f(X+z)$ is exactly stage 11's Husimi symbol at $\hbar=1$, and its
Taylor tensors are the Wiener chaos kernels that v56 grades.
\end{itemize}

\paragraph{Klartag--Ordentlich, SDPI under heat flow.}
\begin{itemize}[leftmargin=*]
\item \textbf{The contraction coefficient.} For $Y=X+\sqrt sZ$, the $\chi^2$ coefficient is the squared maximal correlation $S^2(X,Y)$.
It is bounded between $1-\mathrm{mmse}/\E|X-\E X|^2$ and $(1+s/C_P)^{-1}$, and is at least $e^{-s/C_P}$ for log-concave $X$.
\item \textbf{Log-convexity.} For log-concave $X$, $s\mapsto\log\Var(\E[f(X)|Y])$ is convex.
\item \textbf{Two-flow de Bruijn.} $\frac d{ds}D(\nu_s\|\mu_s)=-\frac12J(\nu_s\|\mu_s)$, with both measures under the flow.
\end{itemize}

\paragraph{Klartag, the evolving ellipsoid.}
\begin{itemize}[leftmargin=*]
\item \textbf{The process.} Start from a ball free of points of a random lattice. Run a Dyson Brownian motion $dA_t=\pi_t(dW_t)$ on
the ellipsoid matrix, where $\pi_t$ projects onto $\{B:\ Bx\cdot x=0$ for every contact point $x\}$. Every lattice point the
ellipsoid touches stays on its boundary.
\item \textbf{The volume accounting.} By It\^o,
\[
d\log\det A_t=\langle A_t^{-1},\pi_tdW_t\rangle-\tfrac12\sum_{i,j}\frac{|\pi_t(u_i\otimes_su_j)|^2}{\lambda_i\lambda_j}\,dt,
\]
and the drift is at most $-\dim F_t/(2\|A_t\|^2)$. Here $\dim F_t\ge n(n+1)/2-|\text{contacts}|/2$.
\item \textbf{The count.} Contacts accrue at rate about $e^{n^2t/8}$ (Siegel's formula and the reflection principle). Maurey's
trick $A_t-a_0=\frac12(W^1+W^2)$ controls $\|A_t\|$. At $T=16n^{-2}\log n$ this gives $\E\log\det A_T\le C-4\log n$: an empty
ellipsoid of volume $cn^2$, hence lattice packings of density $cn^22^{-n}$.
\end{itemize}

\paragraph{Sahasrabudhe, exponentially small scales.}
\begin{itemize}[leftmargin=*]
\item \textbf{The common thread.} A random process is guided toward an exponentially unlikely target. Examples: the Lovett--Meka
walk, a Brownian motion that sticks to every constraint slab it hits and continues inside it; the R\"odl nibble for independent sets;
Klartag's ellipsoid.
\item \textbf{Codegree control.} Campos--Jenssen--Michelen--Sahasrabudhe get $(1-o(1))\frac{d\log d}{2^{d+1}}$ packings and spherical codes
from a nibble on a Poisson-discretized graph. Control of \emph{codegrees}, not triangles, makes the random-like heuristic exact.
\item \textbf{Random symmetric matrices.} Inverse Littlewood--Offord theory and approximate negative correlation give an
$e^{-cn}$ singularity bound for random symmetric $\pm1$ matrices.
\end{itemize}

\paragraph{The caustic-collar billiard paper.} A smooth strictly convex positively curved table whose grazing annulus is
continuously foliated by invariant essential curves is an ellipse. The proof runs in four steps.
\begin{enumerate}[leftmargin=*]
\item \textbf{Graphs and periodic leaves.} Twist and area recurrence make every leaf a graph $p=g_b(\theta)$ in line coordinates, and
make every small rational rotation $l/q$ a wholly periodic leaf.
\item \textbf{Action constancy.} A closed stationary chain on a graph has its action constant in the starting point. Its $\pm q$
Fourier rows $C_q(u)$ therefore vanish. Their derivative is the identity on high modes plus $O(1/q)$, so a Newton scheme on shrinking
complex strips forces the curvature datum $u$ to be analytic.
\item \textbf{Joint analyticity.} A jointly analytic conjugacy $V(y,t)$ is built across the rational resonances. The periodic leaves
supply the compatibility of residual values and of first parameter jets.
\item \textbf{Caustics.} The analytic line graphs are support functions $g(\theta,\lambda)$, with $\lambda=t^2$, $g+g_{\theta\theta}>0$ and
$g_\lambda<0$. Their envelopes $\Gamma_\lambda=gn_\theta+g_\theta e_\theta$ are nested convex caustics filling a collar, and the companion theorem
concludes.
\end{enumerate}

\section{The Schur hub is Eldan's cumulant drift}
\label{sec:hierarchy}
\subsection{Eldan's cumulant hierarchy is the law of total cumulance in time}
Let $\mu$ be a compactly supported probability measure on $\R^d$. Run Eldan's normalized localization: $\mu_t\propto
e^{\langle c_t,x\rangle-\frac12\langle Q_tx,x\rangle}\mu$ with $dc=A_t^{-1}a_tdt+A_t^{-1/2}dB_t$ and $dQ=A_t^{-1}dt$. Write $\kappa_m(t)$ for the
cumulant tensors of $\mu_t$.
\begin{proposition}[Bizeul--Klartag--Lehec, Lemma 4.3; known]\label{prop:bkl}
$d\kappa_m=\kappa_{m+1}(\cdot,\dots,\cdot,A_t^{-1/2}dB_t)-(m\kappa_m+L_m)\,dt$, with $L_m$ as in Section~\ref{sec:readings}. In particular:
\begin{itemize}[leftmargin=*]
\item $L_3=0$, so $\E\kappa_3$ only decays, at rate $3$.
\item The fourth-order term is
\[
L_4(h_1,\dots,h_4)=\sum_{j=2}^4\big\langle\kappa_3(h_1,h_j,\cdot),A^{-1}\kappa_3(h_{\{2,3,4\}\setminus\{j\}},\cdot)\big\rangle .
\]
For instance, $L_4(u,u,u,u)=3\langle\kappa_3(u,u,\cdot),A^{-1}\kappa_3(u,u,\cdot)\rangle$.
\end{itemize}
\end{proposition}
\begin{proposition}[the drift is total cumulance; proved]\label{prop:totalcum}
$\mu_{dt}$ is the posterior of $X\sim\mu$ given one infinitesimal Gaussian observation. Brillinger's law of total cumulance over
set partitions $\pi$,
\[
\kappa_m(X)=\sum_\pi\kappa_{|\pi|}\big(\kappa_{B}(X\,|\,\mathcal G)\big)_{B\in\pi},
\]
applied to that observation, reproduces the drift $-(m\kappa_m+L_m)$ exactly:
\begin{itemize}[leftmargin=*]
\item the one-block partition gives $\E\kappa_m(dt)$;
\item two-block partitions with a singleton give $m\kappa_m\,dt$;
\item two-block partitions with both blocks of size $\ge2$ give $L_m\,dt$;
\item partitions into three or more blocks contribute $o(dt)$.
\end{itemize}
\end{proposition}
\begin{proof}
Conditionally on the observation, $\kappa_k(dt)=\kappa_k+\kappa_{k+1}(\cdot,A^{-1/2}dB)+O(dt)$. The covariance over $dB$ of the fluctuations of
$\kappa_{|B_1|}$ and $\kappa_{|B_2|}$ is $\langle\kappa_{|B_1|+1},A^{-1}\kappa_{|B_2|+1}\rangle dt$. A singleton block $B_1=\{i\}$ has $\kappa_2=A$, and
$\langle A h_i,A^{-1}\kappa_m(h_{\rm rest},\cdot)\rangle=\kappa_m(h)$. Joint cumulants of three or more $O(\sqrt{dt})$ Gaussian-linear fluctuations
vanish at order $dt$.
\end{proof}
So $L_m$ is the \emph{between-path} part of the $m$-th cumulant: what conditioning on the localization path moves out of the
within-path cumulants. It is quadratic in lower cumulants, and the metric is the inverse covariance, because the observation noise is
$A^{-1/2}dB$. The identity is checked on a $7$-point measure in $\R^3$ by Gauss--Hermite averaging over one step (Richardson
in $dt$).
\begin{itemize}[leftmargin=*]
\item The drift of $\kappa_3$ matches $-3\kappa_3$ to $1.5\times10^{-7}$.
\item The drift of $\kappa_4$ matches $-(4\kappa_4+L_4)$ to $2.8\times10^{-7}$; dropping $L_4$ errs by $21\%$.
\item $L_4(u^4)$ equals $3\langle\kappa_3(u,u,\cdot),A^{-1}\kappa_3(u,u,\cdot)\rangle$ to all digits.
\end{itemize}

\subsection{The Schur hub is the between-tier fourth cumulant}
\begin{proposition}[static form; proved]\label{prop:static}
For any $\sigma$-algebra $\mathcal G$ and scalar $X$,
\begin{multline*}
\kappa_4(X)=\E\kappa_4(X|\mathcal G)+4\Cov\!\big(\E[X|\mathcal G],\kappa_3(X|\mathcal G)\big)+3\Var\!\big(\Var(X|\mathcal G)\big)\\
+6\,\kappa\big(\E[X|\mathcal G],\E[X|\mathcal G],\Var(X|\mathcal G)\big)+\kappa_4\big(\E[X|\mathcal G]\big).
\end{multline*}
If $\mathcal G\supseteq\sigma(Y)$ for a vector $Y$ with covariance $C_Y$, then the projection inequality gives
\[
3\Var(\Var(X|\mathcal G))\ \ge\ 3\,\Cov(\Var(X|\mathcal G),Y)\,C_Y^{-1}\,\Cov(Y,\Var(X|\mathcal G)).
\]
Equality holds when the conditional variance is affine in $Y$.
\end{proposition}
Read in the two-tier language of stage 12:
\begin{itemize}[leftmargin=*]
\item the upper tier is the conditioning ($\mathcal G$, the localization centre, i.e.\ the field the next layer sees);
\item the lower tier is the conditional law;
\item the term $3\Var(\Var(X|\mathcal G))$ is the \emph{between-tier} fourth cumulant, and its Schur projection is the
$\kappa_3A^{-1}\kappa_3$ form of $L_4$.
\end{itemize}
The hierarchy's term has the same shape. For centred $X$, $\kappa_3(u,u,\cdot)=\Cov((u\cdot X)^2,X)$, so
\[
L_4(u^{\otimes4})=3\,\Cov((u\cdot X)^2,X)\,A^{-1}\,\Cov(X,(u\cdot X)^2)=3\Var\big(\mathrm{proj}_{\mathrm{span}X}(u\cdot X)^2\big)
\]
under $\mu_t$, averaged over paths: three times the variance of the linear regression of the square on the field. v56's path class is a
regression of the same kind. Note XLI writes it as
\[
P4_i=12\,[YC^{-1}Y^\top]_{ii}=3\Var\big(\mathrm{proj}_{\mathrm{span}\,z_l}\,\mathsf C_i\big),\qquad 2Y_{ij}=\Cov(\mathsf C_i,z_{l,j}),
\]
where $\mathsf C_i$ is the cross term between neuron $i$'s first and second chaos. Using the cross term instead of the whole square removes two
within-tier pieces: the square of the first chaos, which the Gaussian closure already handles, and the square of the second chaos, the
$4$-cycle.

In the chaos-2 model $z=L\cdot x+\frac12(x^\top Hx-\tr H)$, the cross term projects onto the input's linear span as $2(HL)\cdot x$, so
$3\Var$ of it is $12|HL|^2$, the path class of $\kappa_4=12|HL|^2+3\tr H^4$. Under localization this projection is the part of the
conditional variance that moves linearly with the centre. Indeed $\Var(z|\mathcal F_t)=|L+Hm_t|^2/(1+t)+\tr H^2/2(1+t)^2$, whose
linear fluctuation is $2L^\top Hm_t/(1+t)$. So the path class is the between-tier part of the fourth cumulant, read through the field.

\begin{corollary}[what it explains in v56; interpretation, with note XLI's data]
\begin{enumerate}[leftmargin=*]
\item \textbf{The metric is $C^{-1}$, not $\diag(C)^{-1}$.} The observation noise of localization is $A^{-1/2}dB$.
Measured: the hub with the diagonal metric is $30$--$60$ times too large.
\item \textbf{The cumulants must be conditional (quenched).} The drift involves $\E[\kappa_3(t)A_t^{-1}\kappa_3(t)]$ over paths, not
$\E\kappa_3\,A^{-1}\E\kappa_3$. The annealed slice $D21=\kappa(z_i,z_i,\cdot)$ mixes within-tier and between-tier pieces. Measured: the
symmetrized-D21 control removes $0$--$6\%$, the source-resolved hub $64$--$81\%$.
\item \textbf{The weight is fixed.} The three pairings of $L_4(e_i^{\otimes4})$ are the $3$ in $3\Var$. Measured: free-running prefers the
derived weight $1$, and renormalization leaves it at $1.007$.
\item \textbf{Each half alone is adverse.} Proposition~\ref{prop:static} has five terms, and the third-chaos star is the
$4\Cov(\E,\kappa_3)$-type partner. Measured: either piece alone is $+79\%$ to $+131\%$; together they give $-21\%$.
\end{enumerate}
\end{corollary}

\subsection{Predictions: the other pairings from the same solve}
The hierarchy gives every slice of $L_4$, with coefficients:
\begin{gather*}
L_4(e_i,e_i,e_j,e_j)=\langle D_i,A^{-1}D_j\rangle+2\langle\kappa_3(e_i,e_j,\cdot),A^{-1}\kappa_3(e_i,e_j,\cdot)\rangle,\\
L_4(e_i,e_i,e_i,e_j)=3\langle D_i,A^{-1}\kappa_3(e_i,e_j,\cdot)\rangle,
\end{gather*}
with $D_i=\kappa_3(e_i,e_i,\cdot)$. In v56's dressed form the first term of the $(2,2)$ slice is note XLI's ``cheap path part''
$4[YM^{-1}Y^\top]_{ij}$ (one product from the existing solve). The second term and the $(3,1)$ slice need the off-diagonal third cumulant
contracted once: $\sum_k\kappa_3(i,j,k)\,y_{ik}$ with $y_i=C^{-1}D_i$ already formed. On the carried CP sources this is one hub-type
contraction per source-layer, about the cost of the existing D21 hub.
\begin{conjecture}[testable on note XLI's dumps]
The source-resolved $(3,1)$ pairing
\[
3\sum_k\kappa_3(i,j,k)\,[C^{-1}Y]_{ik}
\]
removes a substantial part of the chain's K31 error. That
error is $74\%$ rms; about $40\%$ of the (3,1) slice is ``history'', not a function of the pair state. The pairing uses the full source
tensor, not the pair state, so it is exactly the kind of term the flex theorem allows to carry new information.
\end{conjecture}
\begin{remark}[$\kappa_3$ has no Schur term]
$L_3=0$: under localization the expected third cumulant only decays. A per-layer closure of $\kappa_3$ therefore has no
between-tier correction to add. Its errors are transport errors. This matches note XXXIX: the chain's D21 error is the flat
incoherent sector of the transport (the product-gate defect), right to $2\%$ in its coherent part.
\end{remark}

\section{The critical gain is a conserved dilation charge}
\label{sec:dilation}
Note XLIV (section 7) located the largest structured part of v56's late covariance error. It is one number per layer: the
eigenvalue of the pre-activation covariance along the mean direction, biased low by $0.5$--$0.8\%$. This mode is the dilation
(gain), its error is transported with factor about $0.95$, and its relative variance $\tau$ is roughly stationary at
$0.008$--$0.010$ from layer $6$. Note XLIV left one theory item open: an exact law for how the gain's increments compose through a
\textsc{ReLU} layer. Here it is.

\subsection{The charge}
\begin{proposition}[the dilation charge is conserved; proved]\label{prop:charge}
Let $z=(1+\delta)g$ with $\delta$ independent of $g$, $\E\delta=0$, $\Var\delta=\varepsilon$ and $1+\delta>0$. Then
$\relu(z)=(1+\delta)\relu(g)$, and the next layer is $W\relu(z)=(1+\delta)\,W\relu(g)$. The family of common-scale mixtures is
invariant under the layer map, and the mixing variance $\varepsilon$ is transported unchanged. In covariance form,
$\Cov(z)=(1+\varepsilon)\Cov(g)+\varepsilon\,\E g\,\E g^\top$: the charge lives on the mean direction.
\end{proposition}
This is the exact reason the mode is critical. An error in $\varepsilon$ is carried with factor exactly $1$, while every other direction is
contracted: a layer-$3$ injection loses about $9\times$ in squared norm by layer $15$ (note XLII). Two other views say the same.
\begin{itemize}[leftmargin=*]
\item In stage 12's hyperbolic upper tier, dilation is translation along the kink geodesic, and He initialization is the unique
scale with zero drift along it.
\item In the information-theoretic language of Klartag--Ordentlich (an analogy): a channel contracts every perturbation it does not
preserve, and the layer map preserves the scale-mixture family exactly. The dilation orbit is the one direction with contraction
coefficient $1$.
\end{itemize}

\subsection{The size law}
\begin{proposition}[the dilation sector's leading law; stage 12, Proposition 3.3, at leading order in $1/n$]\label{prop:size}
Let $V_l=\Var_x\log|h_l|^2$ for a fixed network. Its across-input increments are
\[
V_l-V_{l-1}=\frac4n\Big(\frac32-\frac{J_2(\theta_{l-1})}{2\pi}\Big)=\frac{4\theta_{l-1}^2}{n}+O(\theta^4/n),\qquad V_0=\frac2n,
\]
so $nV_l\to44.6$. The relative variance of the mean-direction mode is
\[
\tau_l\simeq\frac{V_l}4+\frac1{nt_l},
\]
where the second term is the bulk's own projection on the mean direction and $t_l$ is the internal clock.
\end{proposition}
At $n=1024$ the leading law gives $V_l/4=0.0051$, $0.0058$, $0.0064$, $0.0068$, $0.0072$ at $l=6,8,10,12,14$, with limit $0.0109$. The input radius
contributes $2/(4n)=4.9\times10^{-4}$, which is note XLIV's ``$1/(2n)$''. Measured values are:
\begin{itemize}[leftmargin=*]
\item on stage 12's network (seed 5), the Rayleigh quotient along the mean is $0.0070$, $0.0077$, $0.0085$ at $l=8,12,16$, with measured
$V_l$ about $7\%$ above the leading law;
\item note XLIV's $\tau$, which uses the top eigenvalue and so sits above the Rayleigh quotient, is $0.0087$--$0.0100$ (network 0) and
$0.0077$--$0.0094$ (network 1) at layers $6$--$14$.
\end{itemize}
So the law has the right order, sign and slow saturation, and is $20$--$30\%$ low at leading order. The quenched value for a given
network is an $O(n^2)$ sum: $\sum_{a,b}\Cov(\relu(z_a)^2,\relu(z_b)^2)$ over the chain's own correlations, a bivariate degree-two
arc-cosine kernel per pair.

\begin{remark}[saturation, not mean reversion; a falsifiable prediction]
Note XLIV read the stationary $\tau$ as a mean-reverting gain. The increment law gives a different mechanism. The gain's across-input
increments are independent to leading order: the weights are fresh at every layer. Their variance $4\theta^2/n$ is summable, because two
inputs at angle $\theta$ disagree on only a fraction $\theta/\pi$ of the gates, so deep layers act alike on all inputs. $\tau$ therefore saturates
without any restoring force.
\begin{itemize}[leftmargin=*]
\item \emph{Test:} the across-input correlation of successive log-norm increments $\Delta_l=\log(|h_l|^2/|h_{l-1}|^2)$.
\item \emph{Prediction:} near $0$, with $\Var\Delta_l\approx4\theta_{l-1}^2/n$.
\item Mean reversion would make the correlation negative.
\end{itemize}
\end{remark}

\subsection{When the charge moves the means}
\begin{theorem}[first-order cancellation; proved, and checked]\label{thm:cancel}
Let $z=(1+\delta)(\mu+s\xi)$ with $\xi\sim N(0,1)$, $\delta$ independent and symmetric, $\Var\delta=\varepsilon$ and $1+\delta>0$, and put $r=\mu/s$. The
Edgeworth-corrected Gaussian closure uses the mixture's own variance $S^2=s^2(1+\varepsilon)+\varepsilon\mu^2$, skewness $\gamma_1$ and excess
kurtosis $\gamma_2$:
\[
S\Big[m_1(R)-\tfrac{\gamma_1}6\He_1(R)\varphi(R)+\tfrac{\gamma_2}{24}\He_2(R)\varphi(R)\Big],\qquad R=\mu/S.
\]
It equals the exact $\E\relu z=s\,m_1(r)$ up to $O(\varepsilon^2)$. To first order the three corrections are
\[
\underbrace{\tfrac\varepsilon2(1+r^2)}_{\text{variance}}\ \underbrace{-\,\varepsilon r^2}_{\gamma_1=6\varepsilon r}\ \underbrace{+\,\tfrac\varepsilon2(r^2-1)}_{\gamma_2=12\varepsilon}\ =0
\qquad(\text{in units of }s\varphi(r)).
\]
\end{theorem}
\begin{proof}
$\kappa_3=6\mu s^2\varepsilon$ and $\kappa_4=12s^4\varepsilon+O(\varepsilon^2)$, by expanding $z-\E z=\delta\mu+s\xi+\delta s\xi$. Then
$\partial_{S^2}[Sm_1(\mu/S)]=\varphi/2S$, $\gamma_1=6\varepsilon r+O(\varepsilon^2)$ and $\gamma_2=12\varepsilon+O(\varepsilon^2)$. Collect the $O(\varepsilon)$ terms.
\end{proof}
Checked for two-point and uniform scale laws at three $(\mu,s)$: the closure error falls by $3.8$--$4.0$ each time $\varepsilon$ halves
($1.1\times10^{-4}\to2.8\times10^{-5}\to7.0\times10^{-6}$ at $\mu=0.7$, $s=1$).

\begin{corollary}[the design rule; derived]\label{cor:package}
The critical mode moves the neuron means only through \emph{inconsistency}. The dilation charge must enter the variance, the third
cumulant and the fourth cumulant as one scale-mixture package. To first order in $\varepsilon$ and in the chain's notation, with
$\widetilde C$ the non-dilation covariance:
\[
\Delta C=\varepsilon(\widetilde C+\mu\mu^\top),\qquad \Delta\kappa_3(a,b,c)=2\varepsilon\big(\mu_a\widetilde C_{bc}+\mu_b\widetilde C_{ac}+\mu_c\widetilde C_{ab}\big),\qquad
\Delta\kappa_4(a,a,a,a)=12\varepsilon\widetilde C_{aa}^2 .
\]
In particular $\Delta D3_a=6\varepsilon\mu_a\widetilde C_{aa}$ and $\Delta D21_{ac}=2\varepsilon(2\mu_a\widetilde C_{ac}+\mu_c\widetilde C_{aa})$. v56 computes these
three pieces by three different routes:
\begin{itemize}[leftmargin=*]
\item the covariance exactly, as $WCW^\top$;
\item $\kappa_3$ through transported legs and births;
\item $\kappa_4$ through the memoryless regeneration.
\end{itemize}
Their dilation parts are therefore mutually inconsistent at the measured $0.5$--$0.8\%$. Note XLIV's injection ledger says the same:
D21 under-feeds the mode at every layer and K31 partly compensates. The rule is to carry $\varepsilon_l$ explicitly and write the three
package terms from $(\varepsilon_l,\mu_l,\widetilde C_l)$ in place of the chain's own dilation components. The cost is $O(n^2)$ per layer.
\end{corollary}
\begin{remark}[the true law is $70\%$ package; measured]
On stage 12's network the truth's marginal cumulants follow the package law in form. The skewness is linear in $r$ ($R^2=0.86$--$0.91$),
and the excess kurtosis is flat. The amplitudes are $\eta=0.71$--$0.72$ (skewness) and $\eta'=0.57$--$0.63$ (kurtosis) of the
full-$V/4$ prediction. The package predicts excess kurtosis $=2\times$ the skewness slope. Measured, the ratio is $1.59$, $1.75$, $1.77$ at
$l=8,12,16$. So the truth contains one approximately consistent scale-mixture component of size $\varepsilon\approx0.005$ at $l=16$, against
$V/4\approx0.008$. The chain can estimate it without truth: regress its own $\gamma_1$ on $r$, set $\varepsilon=\text{slope}/6$, and impose the
package. Deriving $\eta$ is open.
\end{remark}
\begin{remark}[Klartag's principle, applied]
Klartag's ellipsoid and the Lovett--Meka walk keep every constraint they touch and continue in the constraint-preserving subspace.
The dilation package is such a constraint. Along it the chain's transport has no damping, so it must be held exactly rather than
approximated. Every other direction forgets its errors.
\end{remark}

\section{The spherical geometry: support functions, zonoids and mean width}
\label{sec:support}
\subsection{The network is a difference of support functions}
A bias-free \textsc{ReLU} network is positively homogeneous: $f(x)=|x|f(\hat x)$. Under $\gamma$ the radius and the direction are
independent, so $\E f=\E|x|\cdot\int_{S^{n-1}}f\,d\sigma$: the problem lives on the sphere. Every positively homogeneous piecewise-linear
function is a difference of support functions of polytopes (the tropical description of Zhang--Naitzat--Lim). Write
$h_K(x)=\max_{y\in K}\langle y,x\rangle$.
\begin{proposition}[tropical decomposition; proved]\label{prop:tropical}
$\relu(h_A-h_B)=h_{\mathrm{conv}(A\cup B)}-h_B$. Hence, inductively, $f=h_{P_+}-h_{P_-}$ for two polytopes $P_\pm$ built from the rows by
Minkowski sums and convex hulls. A first hidden layer contributes zonotopes $\sum_ac_a[0,w_a]$.
\end{proposition}
\begin{theorem}[the mean is a difference of first intrinsic volumes; proved]\label{thm:V1}
\[
\E f=\frac{\E|x|}2\big(w(P_+)-w(P_-)\big)=\frac{V_1(P_+)-V_1(P_-)}{\sqrt{2\pi}},
\]
where $w$ is the mean width and $V_1$ the first intrinsic volume. Moreover $V_1(P)=\sum_{\text{edges }e}|e|\,\gamma(e,P)$, the edge
length times the external angle. Under the normal-fan duality this is stage 11's border formula,
$\E f=(2\pi)^{-1/2}\sum_F\kappa_F\alpha_F$:
\begin{itemize}[leftmargin=*]
\item the cells of the activation tiling are normal cones of vertices;
\item the walls are normal cones of edges;
\item the kink $\kappa_F$ is the edge length $|e|$;
\item $\alpha_F$ is the external angle.
\end{itemize}
\end{theorem}
\begin{proof}
$\E h_K(x)=\E\sup_{y\in K}\langle y,x\rangle=V_1(K)/\sqrt{2\pi}$ (Sudakov--Tsirelson), and polar coordinates give the mean-width form.
The edge formula is standard (Schneider, Section 4.2).
\end{proof}
Checked in $\R^2$ with $8$ hidden units and a \textsc{ReLU} output.
\begin{itemize}[leftmargin=*]
\item The identity of Proposition~\ref{prop:tropical} holds to $6\times10^{-16}$.
\item The quadrature mean $0.16269286$ equals $\frac{\E|x|}2\cdot\frac{\mathrm{per}(P_+)-\mathrm{per}(P_-)}\pi$ to all printed digits.
\end{itemize}
The packing and convex-geometry literature is therefore natively about our objects. A network is a pair of polytopes, its mean is an
intrinsic volume, its walls are a normal fan, and its gates are a hyperplane tessellation of the sphere. For the last, two inputs at
angle $\theta$ are separated by a fraction $\theta/\pi$ of the hyperplanes (stage 11).

\subsection{Each layer evaluates the support function of a zonoid}
For a law $\mu$ on $\R^n$, the zonoid $Z(\mu)=\{\E[Xg(X)]:\ 0\le g\le1\}$ (Koshevoy--Mosler) has support function $h_{Z(\mu)}(w)=\E\langle w,X\rangle_+$.
\begin{proposition}[the layer map is support-function evaluation; proved]
$\E h_{l,a}=h_{Z(\mu_{l-1})}(w_{l,a})$. The next layer's means are the values of one convex body's support function at the $n$ rows.
The Gaussian closure replaces $Z(\mu_{l-1})$ by the zonoid of $N(\mu,C)$, whose support function is $s\,m_1(r)$. For centred laws this
zonoid is the ellipsoid $C^{1/2}B/\sqrt{2\pi}$. In general it is an ``egg'' interpolating between that ellipsoid and the segment to the mean.
\end{proposition}
\begin{proposition}[the curvature of the zonoid is the wall density; proved]\label{prop:curv}
$\nabla_w^2\,h_{Z(\mu)}(w)=\E[\delta(\langle w,X\rangle)\,X\otimes X]$. In the plane, the radius of curvature $h+h''$ is the density of the projected
law at the wall, weighted by the transverse second moment.
\end{proposition}
Checked for a three-component Gaussian mixture in $\R^2$: the finite-difference Hessian and the wall integral agree to six digits.
This joins three pictures:
\begin{itemize}[leftmargin=*]
\item stage 12's wall-density law (the lower tier's activity is the zonoid's curvature);
\item stage 11's boundary layer;
\item the billiard paper's positivity $g+g_{\theta\theta}>0$.
\end{itemize}
In a $2$-plane, the angular harmonics of $h_Z$ away from the Gaussian egg are the cumulant corrections by order: $\kappa_3$ feeds odd
harmonics and $\kappa_4$ even ones. Non-Gaussianity is non-ellipticity of zonoids.

\subsection{Klartag's ellipsoid in our coordinates (dictionary; analogy, labelled)}
\begin{center}\small
\begin{tabular}{>{\raggedright\arraybackslash}p{5.0cm}>{\raggedright\arraybackslash}p{8.8cm}}\toprule
Klartag's process & the network and the chain\\\midrule
ellipsoid $\{A_tx\cdot x<1\}$, $n(n+1)/2$ degrees of freedom & the closure's Gaussian state $(\mu,C)$; its $n(n+1)/2$ covariance degrees of freedom\\
lattice point on the boundary (contact) & a frozen (resolved) gate, $|r_a|\gg1$; the ball no longer crosses that wall\\
free directions $F_t$, $\dim F_t\ge n(n+1)/2-|\text{contacts}|/2$ & unresolved gates, a fraction $\theta/\pi\approx3/\tau$ per layer (stage 12)\\
drift of $\log\det A_t$, $-\dim F_t/2\|A\|^2$ & the information rate of the gate filtration: log-volume of the input cell lost per layer\\
contacts accrue at rate about $e^{n^2t/8}$ & gates freeze like $\sqrt T$, with $T\approx\theta^2/2$ the internal temperature\\
sticky constraint $\langle A_t,x\otimes x\rangle=1$ & an exact invariant the chain must keep: the dilation package, the Euler--Stein ladder\\\bottomrule
\end{tabular}
\end{center}
The structural content is the sticky principle. A process that is only approximately right gains most by holding exactly the
constraints along which it cannot damp its own errors. In Klartag's case these are the contacts, which keep the ellipsoid lattice-free.
In ours they are the critical (charge) directions of Section~\ref{sec:dilation}.

\paragraph{Spherical codes and codegrees (analogy).} The rows of $W_1$ form a random spherical code, and the first layer's cells are
the $2^n$ images of orthants. In the packing work of Campos, Jenssen, Michelen and Sahasrabudhe, the nibble's ``random-like''
heuristic is made exact by controlling \emph{codegrees}. In ours, the closure (independent, dichotomized-Gaussian gates) is the random-like heuristic, and it is
accurate when pair couplings are small. Stage 12 found the bulk couplings $O(n^{-1/2})$ and incoherent. The one coherent,
uncontrolled ``codegree'' is the Mattis (dilation) component. It is again the charge.

\section{Billiards, packing and exponentially small scales: what transfers}
\label{sec:billiards}
\subsection{The billiard paper, read against our objects}
\begin{enumerate}[leftmargin=*]
\item \textbf{The machinery is ours already.} The paper's last step turns invariant line graphs $p=g(\theta,\lambda)$ into convex bodies
through support functions. Three of its objects are ours.
\begin{itemize}
\item \emph{Contact points.} Its envelope $\Gamma_\lambda=g\,n_\theta+g_\theta e_\theta=\nabla h$ is the contact point of each support line. For a
zonoid, $\nabla_wh_{Z(\mu)}(w)=\E[X\1\{\langle w,X\rangle>0\}]$ is the gate-weighted mean, i.e.\ neuron $w$'s first-chaos vector, which
is the chain's ``arm''.
\item \emph{Convexity.} Its condition $g+g_{\theta\theta}>0$ is Proposition~\ref{prop:curv}: positivity of the wall density.
\item \emph{Nesting.} Its condition $g_\lambda<0$ says the family moves inward. In our setting the zonoids of localized laws are a
martingale family: $\E h_{Z(\mu_t)}=h_{Z(\mu)}$, while each path's zonoid shrinks towards its mean. Cooling plays the role of $\lambda$.
\end{itemize}
\item \textbf{The grazing parameter is quadratic.} For the circle the caustic radius is $R\cos(t/2)=R-\frac R8t^2+\dots$, and $\lambda=t^2$
makes the inward motion first order. The network's internal temperature is $T_l=1-\cos\theta_l\approx\theta_l^2/2$ (stage 12): the
same reparametrization at the same end. In both, the grazing (deep) regime is where the advance angle is small, and the right
coordinate is its square. With $T$ as the parameter, depth is linear in $1/\sqrt T$.
\item \textbf{Rigidity has a probabilistic twin.}
\begin{proposition}[proved; elementary]
If every exponential tilt of $\mu$ has the same covariance, $\mu$ is Gaussian. Indeed $\nabla^2\Lambda_\mu$ is then constant, so $\Lambda_\mu$ is
quadratic. Equivalently, if the Gaussian closure is exact along a full open family of tilts (a ``foliation'', not a discrete set),
then all cumulants of order $\ge3$ vanish.
\end{proposition}
The analogy is in the hypothesis. The paper needs a \emph{full} foliation, because a Cantor family of caustics (Lazutkin) exists for
non-ellipses. A closure that is exact on isolated states says nothing; exactness on a continuum forces Gaussianity. The first
obstruction is $\kappa_3$, the first odd harmonic of the zonoid.
\item \textbf{The method: family invariants pin the high modes.} Sliding periodic families have constant action, which gives one exact
equation $C_q(u)=0$ per high frequency, with derivative equal to the identity plus $O(1/q)$ on the high modes. A Newton scheme on
shrinking strips then determines the high-frequency data from the low. Our exact invariants are few:
\begin{itemize}
\item the dilation charge (one direction per layer);
\item the Euler--Stein ladder, $(1-r)f_r=(r+1)(r+2)\tr f_{r+2}$ (one trace per chaos order);
\item the localization variance decomposition $\E A_t+\E a_t\otimes a_t=\Cov$ (one matrix identity per time, automatically satisfied by the
closure).
\end{itemize}
None is diagonal on the flat (incoherent) sector. This is the flex theorem of note XXXVIII seen from the billiard side: the pair state
plus our invariants does not determine the omitted fourth-order classes. What did play the role of the missing equations in v56 was
extra data, the source-resolved legs, read through the Schur hub (Section~\ref{sec:hierarchy}).
\end{enumerate}
\emph{Verdict.} The billiard paper supplies a dictionary (support functions, envelopes, positivity, the quadratic grazing
parameter), a rigidity principle (only full families constrain), and a method. It supplies no algorithm by itself. The parabolic
depth map has no resonances: in Fatou time it is a translation, so the small divisors that drive the billiard proof do not occur in the
chain's depth recursion.

\subsection{Packing and exponentially small scales}
\begin{itemize}[leftmargin=*]
\item \textbf{Why sampling cannot see cells.} The mean is a sum over exponentially many cells, each of Gaussian mass about $2^{-n}$ at
the first layer. At flopscope's budget a Monte Carlo pass affords about $6\times10^4$ samples. That gives per-neuron MSE of order
$10^{-6}$--$10^{-5}$, against the chain's $1.5\times10^{-8}$. The ``guide the process to the rare event'' lesson therefore transfers only in its
deterministic form: guided, sticky evolution of the closure's state.
\item \textbf{Codegrees.} The Campos--Jenssen--Michelen--Sahasrabudhe insight, that controlling codegrees makes a random-like
heuristic exact, has the reading of Section~\ref{sec:support}: the closure is accurate when the gate gas's pair couplings are small,
and the one coherent coupling is the dilation.
\item \textbf{Anti-concentration.} The Littlewood--Offord theme, small-ball probabilities of $\sum_aw_ah_a$, is the wall density of
stage 12. Its law $(\int Q)\sin(\theta/2)/\sqrt\pi$ is an anti-concentration statement for the network's own fields.
\end{itemize}

\section{Heat flow, slicing, and a memory budget}
\label{sec:heat}
\subsection{A memory budget: negentropy as an integrated localization deficit}
\begin{proposition}[Klartag--Lehec, Corollary 3.2 with Lemma 3.3; restated]\label{prop:negent}
For an isotropic law $\mu$ on $\R^n$, with $A_t$ its covariance process under Eldan's (standard) localization,
\[
D(\mu\,\|\,\gamma)=\tfrac12\int_0^\infty\Big(\frac n{1+t}-\E\Tr A_t\Big)dt .
\]
For the Gaussian $A_t=\mathrm{Id}/(1+t)$ exactly, so the integrand is the localization deficit: how much faster than a Gaussian the
law localizes.
\end{proposition}
\begin{proof}
$D(\mu\|\gamma)=\frac12\int_0^1\E|v_r|^2dr$ with $\E v_r\otimes v_r=(\mathrm{Id}-\E\Gamma_r)/(1-r)$, $\Gamma_r=(1+t)A_t$ and $r=t/(1+t)$. Change variables.
\end{proof}
Checked in one dimension by Monte Carlo over the localization path: logistic law, $0.0144$ against $0.0142$; two-bump mixture,
$0.1720$ against $0.1728$.

\emph{Why this matters here.} Stage 12 asked for a modular quantity that controls the closure's error, and named the negentropy
of $z_l$ relative to its quasi-free projection as the candidate. Proposition~\ref{prop:negent} turns that negentropy into a time integral of
second moments of \emph{localized} laws. By stage 12, the network's internal clock is a localization time. The chain could therefore
estimate its own memory budget from quantities of the same kind it carries: the conditional (quenched) covariances along the
localization, the same objects the Schur hub reads. Whether the budget predicts the per-layer closure error is a cheap, truth-free
test.

\subsection{Shannon--Stam stability: the quenched content as covariance fluctuation}
Eldan--Mikulincer (as used by Klartag--Lehec) bound $\int\E|\Gamma_r-\E\Gamma_r|^2\,dr/(1-r)$ by the entropy gain of $\mu$ under
$X\mapsto(X+X')/\sqrt2$. That integral is the variance, across localization paths, of the localized covariance, which is exactly the
quenched content that separates the Schur hub from its annealed control (Section~\ref{sec:hierarchy}). Its volatility is the third cumulant:
$dA_t=\sum_iH_{i,t}dB_{t,i}-A_tdt$ with $H_{i,t}=\kappa_3^{\mu_t}(\cdot,\cdot,A_t^{-1/2}e_i)$ (Bizeul--Klartag--Lehec, eq.~(69)). The closure's
ellipsoid is deterministic. The true ellipsoid diffuses, with $\kappa_3$ as its diffusion coefficient. This is Klartag's
stochastically evolving ellipsoid with the lattice replaced by the law's own cumulants.

\subsection{Heat-flow contraction and the chaos spectrum}
For the Gaussian input, the posterior mean of the output under the observation $X+\sqrt sZ$ is $\E[f(X)|X+\sqrt sZ]$, the Husimi symbol
at the posterior centre. Its variance is
\[
\Var\,\E[f(X)\,|\,X+\sqrt sZ]=\sum_{k\ge1}(1+s)^{-k}\,\|f_k\|^2,
\]
with $f_k$ the $k$-th Wiener chaos component. This is log-convex in $s$, as Klartag--Ordentlich prove for every log-concave input. The
chaos spectrum that v56 grades is thus the sequence of power-series coefficients, in $u=1/(1+s)$, of one measurable curve: the
posterior-mean variance against noise level. The maximal correlation of the Gaussian channel is $(1+s)^{-1/2}$, attained by the first chaos. For the network output, the ratio
$\Var\E[f|Y]/\Var f$ falls below $(1+s)^{-1}$ by exactly the share of the higher chaos. Note XLII measured the first-chaos share of
deep pre-activations falling with depth, consistent with v56 needing the second and third chaos. One more link: the tilt averages of Bizeul--Klartag--Lehec
are, for the Gaussian, $F_f(z)=\E f(X+z)$. Their Taylor tensors are the chaos kernels, and the Song--Zhang criterion turns geometric
decay of those kernels into a Poincar\'e bound. For the chain, the same decay rate controls the error of truncating the chaos grading.

\section{What to build next, and what to test}
\label{sec:next}
\subsection{Builds, ranked by expected value per unit cost}
Each item is stated so that it can be pre-registered in the protocol of note XXXIX: cold screen on networks $0$--$15$, then scored
with refitted counterterms and judged held-out.
\begin{enumerate}[leftmargin=*]
\item \textbf{The dilation-consistent closure} (Section~\ref{sec:dilation}).
\begin{itemize}
\item \emph{Build.} Carry $\varepsilon_l$. Estimate it from the chain's own D3 slope, or from the $O(n^2)$ degree-two arc-cosine sum.
Overwrite the dilation components of C, D3, D21 and the $\kappa_4$ diagonal by the package of Corollary~\ref{cor:package}.
\item \emph{Cost.} $O(n^2)$ per layer.
\item \emph{Mechanism.} Theorem~\ref{thm:cancel}: once the package is consistent, the critical mode no longer moves the means at first
order, however large its accumulated error.
\item \emph{Prediction.} The dilation-mode share of the late covariance error (note XLIV's $c_1$, $0.33\to0.58$) falls substantially, and
raw MSE falls by a fraction of the dilation oracle's value.
\item \emph{Bound.} The oracle (true eigenvalue along the true mean at every layer) should be run first, as note XLIV proposed.
\end{itemize}
\item \textbf{The remaining $L_4$ pairings from the existing solve} (Section~\ref{sec:hierarchy}).
\begin{itemize}
\item \emph{Build.} The $(2,2)$ path part $\langle D_i,C^{-1}D_j\rangle$ is one product from the solve v56 already does. The off-diagonal
pairings $2\langle\kappa_3(ij\cdot),C^{-1}\kappa_3(ij\cdot)\rangle$ and $3\langle D_i,C^{-1}\kappa_3(ij\cdot)\rangle$ need one hub-type contraction
of each source with $y_i=C^{-1}Y_i$.
\item \emph{Prediction.} The $(3,1)$ pairing removes a substantial part of the K31 error ($74\%$ rms), and its free amplitude refits
near $1$.
\item \emph{Order.} First run the diagnostic: regress the truth's K31 on the pairing, on note XLI's dumps, before any chain run.
\end{itemize}
\item \textbf{Sticky truth-free constraints.} After each layer, project the chain's state onto the affine set where the exact identities
hold, by the minimal Fisher--Rao change: the dilation package, ladder rung $0$ ($\E y=2\tr f_2$, unit by unit) and rung $1$ ($\tr f_3=0$).
The cost is $O(n^2)$. The rung-$0$ residual alone is a free measurement of what the legs omit (note XLIV, section 8).
\item \textbf{The memory budget} (Proposition~\ref{prop:negent}). Estimate the localization deficit from the carried conditional covariances,
and test whether it predicts the per-layer injection. This is a diagnostic, not a build.
\end{enumerate}

\subsection{Predictions (falsifiable)}
\begin{enumerate}[leftmargin=*]
\item Successive across-input log-norm increments are uncorrelated, with $\Var\Delta_l\approx\frac4n(\frac32-\frac{J_2(\theta_{l-1})}{2\pi})$. The gain
saturates; it does not mean-revert.
\item The truth's marginal cumulants satisfy, approximately, kurtosis $\approx2\times$ the skewness slope (measured $1.6$--$1.8$). The chain's
own satisfy it less well, and the gap tracks the dilation injection.
\item Imposing the dilation package leaves the counterterms' dilation-sensitive amplitudes (var, D21, K31 at depth) near their
unrenormalized values. The package should absorb what they now compensate.
\item The KLS pairings are law terms: their free amplitudes refit at their derived values ($1$, $2$, $3$ for the three slices).
\end{enumerate}

\subsection{Open problems}
\begin{enumerate}[leftmargin=*]
\item \textbf{Derive the efficiency $\eta\approx0.71$.} This is the fraction of the across-input log-norm variance that acts as a common
scale on the next layer's bulk. It is a four-point function of the arc-cosine hierarchy.
\item \textbf{Integrate the hierarchy along the network's internal clock.} Stage 12 makes each layer a localization step of size
$\Delta\Fa=1$. Integrating $-(4\kappa_4+L_4)$ over the step with the network's own conditional cumulants should give the exact per-layer
$\kappa_4$ update. Its leading term would be the Schur hub, and the correction terms would come with their weights.
\item \textbf{Intrinsic volumes along depth.} $\E f=(V_1(P_+)-V_1(P_-))/\sqrt{2\pi}$. Does the Gaussian closure correspond to a
Steiner-type approximation of the tropical polytopes, and is its error a higher intrinsic volume?
\item \textbf{A rigidity statement for the chain.} An analogue of the caustic-collar theorem: if a closure is exact on a full
neighbourhood of the network's localized states, is the network linear on the support?
\end{enumerate}

\section*{Verification}
The checks are run by \texttt{scripts/verify\_stage13.py}; its output is saved next to this note. The network measurements quoted
from stage 12 come from \texttt{scripts/verify\_stage12.py}. The v56 figures come from notes XXXIX, XLI and XLIV on branch
\texttt{claude/bold-albattani-xgecxi}.
\begin{center}\small
\begin{tabular}{>{\raggedright\arraybackslash}p{5.0cm}>{\raggedright\arraybackslash}p{9.8cm}}\toprule
identity & result\\\midrule
1. Eldan cumulant drift (Prop.~\ref{prop:bkl}, \ref{prop:totalcum}) & $7$-point measure in $\R^3$, Gauss--Hermite over one step, Richardson in $dt$: $\kappa_3$ drift $=-3\kappa_3$ to
$1.5\times10^{-7}$; $\kappa_4$ drift $=-(4\kappa_4+L_4)$ to $2.8\times10^{-7}$ ($L_4$ is $20\%$; without it the error is $21\%$); $L_4(u^4)=3\langle\kappa_3(uu\cdot),A^{-1}\kappa_3(uu\cdot)\rangle$ exactly\\
2. First-order cancellation (Thm~\ref{thm:cancel}) & two-point and uniform scale laws at three $(\mu,s)$: the error falls by $3.8$--$4.0$ per halving of $\varepsilon$, e.g.\
$1.10\times10^{-4}$, $2.77\times10^{-5}$, $6.97\times10^{-6}$\\
3. Dilation size law (Prop.~\ref{prop:size}) & $V_l/4=0.0051$, $0.0058$, $0.0064$, $0.0068$, $0.0072$, $0.0075$ at $l=6,8,10,12,14,16$; input radius $4.9\times10^{-4}$; per-layer
increments $5.4$, $2.9$, $1.8\times10^{-4}$ at $l=6,10,14$\\
4. Tropical mean width (Prop.~\ref{prop:tropical}, Thm~\ref{thm:V1}) & $\R^2$, $8$ units: $\max|\relu(h_A-h_B)-(h_{\mathrm{conv}(A\cup B)}-h_B)|=6\times10^{-16}$; $\E f=0.16269286$ by quadrature and by perimeters\\
5. Zonoid curvature (Prop.~\ref{prop:curv}) & three-component mixture in $\R^2$: Hessian $(0.210493,0.140329,0.093552)$ by finite differences and by the wall integral\\
6. Negentropy budget (Prop.~\ref{prop:negent}) & logistic $0.0144$ against $0.0142$; two-bump $0.1720$ against $0.1728$ (Monte Carlo over paths, $t\le3000$)\\\bottomrule
\end{tabular}
\end{center}

\begin{thebibliography}{99}\small
\bibitem{stage12} Stage 12 note: \emph{the network as its own two-tier system} (this repository, \texttt{notes/stage12}).
\bibitem{stage11} Stage 11 note: \emph{tilings, Bratteli diagrams and singular foliations} (\texttt{notes/stage11}).
\bibitem{XLI} Working note XLI, \emph{chaos grading and the Schur hub}, and note XXXIX, \emph{the leaderboard system} (branch \texttt{claude/bold-albattani-xgecxi}, \texttt{notes/chaos-grading}, \texttt{notes/leaderboard-system}).
\bibitem{XLIV} Working note XLIV, \emph{the fresh-weight theorem}, sections 7--8 (same branch, \texttt{notes/fresh-weight}).
\bibitem{XV} Working note XV, \emph{every computation of the chain} (same branch, \texttt{notes/chain-anatomy}); the base chain is 504aldo's open-source factorized third-cumulant chain, v29 (MIT).
\bibitem{KL} B.~Klartag, J.~Lehec, Affirmative resolution of Bourgain's slicing problem using Guan's bound, arXiv:2412.15044.
\bibitem{BKL} P.~Bizeul, B.~Klartag, J.~Lehec, Presenting a proof of the Kannan--Lov\'asz--Simonovits conjecture, arXiv:2610.05474.
\bibitem{KO} B.~Klartag, O.~Ordentlich, The strong data processing inequality under the heat flow, arXiv:2406.03427.
\bibitem{Kl} B.~Klartag, Lattice packing of spheres in high dimensions using a stochastically evolving ellipsoid, arXiv:2504.05042.
\bibitem{Sah} J.~Sahasrabudhe, Probabilistic combinatorics at exponentially small scales, arXiv:2512.15077.
\bibitem{bill} OpenAI, Continuous phase foliations create analytic caustic collars (24 September 2026), with its companion \emph{Rigidity of smooth billiards with a continuous caustic collar}.
\bibitem{Eldan} R.~Eldan, Thin shell implies spectral gap up to polylog via a stochastic localization scheme, GAFA 23 (2013) 532--569.
\bibitem{EM} R.~Eldan, D.~Mikulincer, Stability of the Shannon--Stam inequality via the F\"ollmer process, PTRF 177 (2020) 891--922.
\bibitem{LM} S.~Lovett, R.~Meka, Constructive discrepancy minimization by walking on the edges, SIAM J.~Comput.~44 (2015) 1573--1582.
\bibitem{CJMS} M.~Campos, M.~Jenssen, M.~Michelen, J.~Sahasrabudhe, A new lower bound for sphere packing, arXiv:2312.10026.
\bibitem{Brill} D.~R.~Brillinger, The calculation of cumulants via conditioning, Ann.~Inst.~Statist.~Math.~21 (1969) 215--218.
\bibitem{ZNL} L.~Zhang, G.~Naitzat, L.-H.~Lim, Tropical geometry of deep neural networks, ICML 2018.
\bibitem{KM} G.~A.~Koshevoy, K.~Mosler, Lift zonoids, random convex hulls and the variability of random vectors, Bernoulli 4 (1998) 377--399.
\bibitem{Schn} R.~Schneider, \emph{Convex Bodies: the Brunn--Minkowski Theory}, 2nd ed., Cambridge 2014 (support functions, intrinsic volumes, mean width).
\bibitem{HN} B.~Hanin, M.~Nica, Products of many large random matrices and gradients in deep neural networks, CMP 376 (2020) 287--322.
\bibitem{CS} Y.~Cho, L.~K.~Saul, Kernel methods for deep learning, NeurIPS 2009.
\bibitem{Clark} C.~E.~Clark, The greatest of a finite set of random variables, Oper.~Res.~9 (1961) 145--162.
\end{thebibliography}
\end{document}
